\subsection{Configuration matrix and predictions}

We combine the two read concerns (\texttt{local}, \texttt{majority}), two write
concerns (\texttt{w:1}, \texttt{majority}), and causal sessions switched off or
on. This produces eight configurations. For each property, we keep the
operation sequence, logical document, and replica-set topology fixed while
changing these controls. Table~\ref{tab:configurations} marks the
configuration-property pairs for which we made a positive prediction; the
other cells are descriptive comparisons, not claims of a guarantee.

\begin{table}[ht]
\centering
\small
\setlength{\tabcolsep}{3pt}
\caption{C1-C8 settings and properties expected under the target schedules.}
\label{tab:configurations}
\begin{tabular}{@{}l l l l p{0.30\linewidth}@{}}
\toprule
Config. & Read concern & Write concern & Causal session & Expected properties \\
\midrule
C1 & local & \texttt{w:1} & off & Descriptive only \\
C2 & local & \texttt{w:1} & on & Descriptive only \\
C3 & majority & \texttt{w:1} & on & MR, WFR \\
C4 & local & majority & on & MW \\
C5 & majority & majority & off & Descriptive only \\
C6 & majority & majority & on & RYW, MR, MW, WFR \\
C7 & majority & \texttt{w:1} & off & Descriptive only \\
C8 & local & majority & off & Descriptive only \\
\bottomrule
\end{tabular}
\end{table}

C1 leaves reads local, acknowledges one replica, and carries no causal
dependency; we therefore expect a stale or reordered result to be possible
when the schedule splits the replicas. C3 asks whether majority reads and a
causal session preserve MR and WFR in the tested sequence even though its
write uses \texttt{w:1}; this is a schedule-specific prediction, not a general
MongoDB guarantee for that combination. C4 tests whether an acknowledged
majority write preserves MW when a successor write follows it; its local read
does not protect WFR. C6 combines majority reads, majority writes, and a
causal session. We expect its dependencies not to produce a completed
counterexample, although a partition may prevent an operation from finishing.

We compare these predictions only with completed, decidable histories. A
timeout or missing precondition leaves the prediction untested in that trial.
In particular, a stale RYW read in C7 describes an unprotected setting; it
does not falsify a guarantee we assigned to C7.

\subsection{Fault expectations}

The first campaign changes settings while its schedule creates controlled
replica divergence. The second asks how the same client dependencies behave
when we change the topology.

\begin{description}
\item[F1, secondary loss] Removing one secondary should leave the primary and
a majority side intact. We expect the scheduled calls to complete without a
new consistency counterexample.
\item[F2, primary loss] Stopping the primary should trigger an election. We
expect calls made after the election barrier to preserve the tested
dependencies, but we make no prediction about a call during the election
because our schedule does not issue one.
\item[F3, replication partition] Isolating the former primary should separate
a \texttt{w:1} acknowledgement from a majority acknowledgement. We predict
that weak settings may complete with stale or reordered state, whereas C6 may
wait or leave a dependent operation unresolved.
\end{description}
